\documentclass[letterpaper,12pt]{article}
\usepackage{helvet}
\usepackage{fancyheadings}
\usepackage{verbatim}
\usepackage{hyperref}
\pagestyle{fancy}
\usepackage{graphicx}
\setlength\textwidth{6.5in}
\setlength\textheight{8.5in}
\begin{document}
\title{Literate Programming in the Era of LLMs}
\author{Ramesh Subramonian }
\maketitle
\thispagestyle{fancy}
\lhead{}
\chead{}
\rhead{}
\lfoot{}
\cfoot{}
\rfoot{{\small \thepage}}


\newtheorem{problem_statement}{Problem Statement}
\newtheorem{notation}{Notation}
\newtheorem{convention}{Convention}
\newtheorem{assumption}{Assumption}

\newcommand{\TBC}{\framebox{\textbf{TO BE COMPLETED}}}
\newcommand{\be}{\begin{enumerate}}
\newcommand{\ee}{\end{enumerate}}
\newcommand{\bi}{\begin{itemize}}
\newcommand{\ei}{\end{itemize}}
\newcommand{\bd}{\begin{description}}
\newcommand{\ed}{\end{description}}

\newcommand{\True}{{\tt true} }
\newcommand{\False}{{\tt false} }
\newcommand{\uint}{{\tt uint32\_t} }
\newcommand{\qt}{{\tt qtype\_t} }
\newcommand{\jt}{{\tt jtype\_t} }
\newcommand{\qdf}{{\tt qdf\_rec\_type} }

\newcommand{\jarr}{{\tt j\_array} }
\newcommand{\jobj}{{\tt j\_object} }

\newcommand{\getq}{{\tt get\_qtype}}
\newcommand{\getj}{{\tt get\_jtype}}
\newcommand{\getn}{{\tt get\_length}}
\newcommand{\getrv}{{\tt get\_read\_arr}}
\newcommand{\getwv}{{\tt get\_write\_arr}}
\newcommand{\getnk}{{\tt get\_num\_keys}}
\newcommand{\getrnn}{{\tt get\_read\_nulls}}
\newcommand{\getwnn}{{\tt get\_write\_nulls}}
\newcommand{\chk}{{\tt chk\_qdf}}

\newcommand{\setq}{{\tt set\_qtype}}
\newcommand{\setj}{{\tt set\_jtype}}
\newcommand{\setn}{{\tt set\_length}}

\newcommand{\mkUA}{{\tt make\_uniform\_array}}
\newcommand{\mknnUA}{{\tt make\_uniform\_array\_with\_nulls}}

\newcommand{\setv}{{\tt set\_arr}}
\newcommand{\setnn}{{\tt set\_nn}}

\newcommand{\Q}{{\tt Q}}
\newcommand{\QDF}{{\tt QDF}}


\abstract{LLMs have taken the programming world by storm. 
How does this change the role of the software engineer? 
Assuming that the cost of writing code will approach zero,
what does a software engineer do? 
We posit that there are a few fundamental changes in the way we will 
build software. (1) We will focus more on writing 
unambiguous, semi-formal specifications than we have traditionally done.
(2) We will deploy more custom code, tailored to the specific need of the
application, thereby reducing code bloat. 
(3) We will adopt just-in-time manufacturing approach to software deployment 
In this paper, we will use an example of moderate complexity to 
demonstrate how these principles are put into practice. 

\section{Introduction}


Traditionally, the cost of writing software 
(and modifying it when new needs arose) was so high that the design of 
code optimized for maintainability. Most seasoned software engineers will have
experienced the sinking feeling when asked to add new requirements to a large
code base, the fear of not fully understanding whether a fix here 
would cause a blow out there. We would argue that assuming 
zero cost of (bug-free) software development flips this paradigm on 
its head. This moves the burden to making sure that the specifications are
unambiguous and consistent and meet the business requirement precisely.

Good software engineering was never about churning out lines of code. Instead,
it was about designing systems that met the requirements 

The answer is to be found in Djsktra's claim that
``Computer science is no more about computers than astronomy is about
telescopes''. Writing efficient code was a means to an end.
There have been a series of improvements in that process, 
LLMs representing the
latest and admittedly most unexpected quantum jump in efficiency.
Writing formal specifications is a declaration of intent. 
Until the singularity arrives, specifications represent the safest 
bet that 

The author makes no claim for what AI will do in the longer term.

\begin{center}
{\em How much code should a coder code if Claude Code could code code?}
\end{center}

\section{Motivation}

Earlier, 
the author had built a column-oriented analytical database called \Q. 
As with most such databases, greater efficiencies are possible
when data is created and accessed within the database
engine itself. 
However, a database that doesn't share its data with other applications isn't of
much use. This necessitated the development of \QDF\ (\Q's data format) for data
exchange, not for bulk import/export.

We considered several other formats, starting with the usual suspects,
 CSV and JSON. For efficiency reasons, we narrowed our focus to binary formats, like BSON
 and AVRO. Ultimately, we decided it best to roll our own format based on the
 following requirements and usage patterns
 \be
 \item ease of adding fixed-width custom types
 \item optimizing computations on the data without moving or
   transforming it
 \item allowing for in-situ operations,
   although most operations are write-once.
 \item optimizing for the case where data is read more 
   often than it is written. 
 \item ability to control whether a \QDF\ packet is ``write-once'' or not
   \ee

The development work could be broadly divided into two parts
\be
\item Definition of the data format and writing helper functions which allowed
  the creation and modification of data packets in \QDF\ format and access to
  elements of such a packet
\item Definition and development of operators that read from (and occasionally
  wrote to) the data packet.
  A simple example of an operator is a dot product which
accesses two numeric arrays of equal length and returns a number. 
  \ee
Over a period of four years, we found that the changes in
the underlying data format were both infrequent and minor.
As a consequence, the work on getters and setters was minimal. 
Instead, most of the development work involved
\be
\item writing new operators. 
\item optimizing the performance of existing operators 
  \ee
This led to the following questions
\be
\item Could the use of LLMs have reduced the development and maintenance cost?
\item If so, what would the software architecture of such a system look like?
  \ee

Giving GenAI the benefit of the doubt, we assumed the answer to the first
question would be Yes and proceeded to answer the second question. 

\section{Designing for LLMs}

\subsection{Choice of Language}

We chose to build our system using C and Lua and \LaTeX.  
There is nothing sacrosanct about these choices. Other 
choices of compiled language and interpreted
language and document formatting language are likely to work just as well. 
The reasons for this choice were
\be
\item The Ousterhout dichotomy. This is a principle of software design that I
  have found very useful. For those unfamiliar with the idea, 
``Ousterhout's dichotomy claims that there are two general categories of programming languages:
\be
\item low-level, statically-typed systems languages
\item high-level, dynamic scripting languages
  \ee
Systems languages are best for efficiently handling large quantities of data, implementing algorithms, and can handle significant internal complexity. Scripting languages are best for glueing other programs together, exploring a problem, and extending applications.
''
\footnote{\url{https://gist.github.com/nat-418/bdccb2428cde750c9f42bf9bbc1a50d3}}

\item Completeness. 
It was very much our intent to
have a fully functional system of moderate complexity which would allow us to
suss out any gotchas that a ``Hello World'' program might have glided by.
The author is sufficiently conversant in these languages to meet this goal.

\item Ease of exposition. 
  The simple syntax of C and Lua make the code snippets widely
  readable
\item Ease of implementation.
``Lua was designed since its inception to interoperate with
other languages, both by extending --- allowing Lua code to call functions
written in a foreign language --- and by embedding --- allowing foreign code to call functions written in Lua. Lua is
thus implemented not as a standalone
program but as a library with a C API.
This library exports functions that cre ate a new Lua state, load code into a
state, call functions loaded into a state,
access global variables in a state, and
perform other basic tasks. The stand alone Lua interpreter is a tiny application
written on top of the library.'' \cite{Lua2019}
\item \LaTeX, when combined with {\tt make}, allows us to ``compile''
  documentation from disparate sources. It allows us to create 
  different targets from common sources. 
\ee



\subsection{Control the Data}
We took our inspiration from the quote attributed to Linus Torvalds --- 
{\em ``Bad programmers worry
about the code. Good programmers worry about data structures and their
relationships.} Hence, we focused our efforts on designing an efficient data
format  (Section~\ref{data_format_defn}) and 
and defining how it would be accessed (Section~\ref{accessor_defn}).

\subsubsection{Define the Data Format}
\label{data_format_defn}

We wanted the data format to resemble JSON but to have efficient 
representations of commonly used types e.g., dataframes. Consider the JSON type
{\tt number}. This is stored in QDF as a 16 byte packet where
\be
\item Byte 0 contains JSON type information --- {\tt number} in this example.
\item Byte 1 contains data type information e.g., {\tt int32\_t} or {\tt double}
\item Byte 2 is the successor offset. While this is usually 0, it helps tell us
  where the successor QDF, if any, is located. This is used when we want to
  make sure that certain data structures are memory aligned. 
\item Byte 3 is unused.
\item Bytes 4 to 7 contain an offset to the parent packet, if any. 
  For example,
  in the JSON array \verb+[ true, "abc", 123 ]+, the number 123 needs to 
  know the start of the parent array of which it is a part.
\item Bytes 8 to 15 contain the actual value
  \ee
\subsubsection{Define the Accessors}
\label{accessor_defn}

After defining the data format, we implemented a number of helper 
functions which broadly fell into these categories
\be
\item readers or getters e.g., calling {\tt get\_jtype()} on a QDF packet tells
  us what JSON type it is i.e., {\tt nil, number, string, \ldots}
\item writers or setters
\item makers or constructors e.g., calling {\tt make\_number()} with the 
  type and value of the number to be created as arguments, 
\item destructors were largely handled by Lua's garbage collection
\item checkers, which validated the syntactic structure of a data packet.
  These can be thought of as invariants that could be asserted before and after
  an accessor was called. This was inspired by \cite{hoare1969axiomatic}.

  \ee

\subsubsection{Test the Accessors}
Once the accessors were in place, it was necessary to write unit tests which
would give us confidence in their implementation. We started by hand writing
these tests but quickly transitioned as described in
Section~\ref{llm_uint_tests}.

\subsubsection{Define the Operators}

We took our inspiration from Knuth: ``Let us change our traditional attitude to the construction of programs: 
Instead of imagining that our main task is to 
instruct a computer what to do, let us concentrate rather on explaining
to human beings what we want a computer to do.'' \cite{knuth1984literate}.

Taking this last sentence to its logical conclusion, we chose to 
concentrate {\bf entirely} on the human explanation, 
leaving the job of instructing the computer what to do to the LLM.

In this paper, Knuth argues that the system he proposes, WEB, is 
itself is chiefly a combination of two other languages: (1) a document formatting language and (2) a
programming language. My prototype WEB system uses
\TeX\ as the document formatting language and PASCAL
as the programming language, but the same prin ciples would apply equally well if other languages were
substituted. Instead of \TeX, one could use a language
like Scribe or Troff; instead of PASCAL, one could use
ADA, ALGOL, LISP, COBOL, FORTRAN, APL, C, etc.,
or even assembly language. The main point is that WEB
is {\bf inherently bilingual}, and that such a combination of
languages proves to be much more powerful than either
single language by itself. WEB does not make the other
languages obsolete; on the contrary, it enhances them.

In our case, we chose to use \LaTeX\ for specifications 
and have the LLM write the code in Lua, C and ISPC \cite{ISPC}
The author's familiarity with these languages allowed for code review.
If LLMs live up to the trajectory claimed for them, 
we expect to gain confidence in them and reduce the need for oversight.

\section{Making it concrete}

We describe the implementation of the operator {\tt coalesce}
\footnote{url{https://www.rdocumentation.org/packages/dplyr/versions/1.0.10/topics/coalesce}}.

\subsection{Starting Out}
\begin{verbatim}
lQDF = require 'lQDF'
\end{verbatim}
This informs the system about the data definitions and data accessors.
To those familiar with LuaJIT \cite{LuaJIT}, we have
\be
\item {\tt ffi.cdef}'d the function declarations
\item {\tt ffi.load}'d the library that implements the accessors
  \ee

  \subsection{Registering an operator}
\begin{verbatim}
local C = require 'coalesce'
lQDF.Q.register('coalesce', C)
\end{verbatim}

{\tt C} is a Lua table that contains the specification of the operator. 
Details are in Section~\ref{coalesce_spec}. It consists of the following
\be
\item \verb+generic_tex_spec_file+: \LaTeX\ file containing specification 
  See Appendix for an example %% TODO 
\item \verb+get_subs+: a function which produces substitutions needed
  to specialize the above \LaTeX\ file based on the types of the arguments
\item \verb+lua_calling_spec+: table containing specification used to 
  create a function {\tt
  run()} which is used at runtime to create (if necesssary) 
  and execute a C function
  \ee

The call to {\tt Q.register} causes the following
\be
\item Records the information provided
\item Uses \verb+lua_calling_spec+ to create the code for function {\tt run()}.
  Currently, this is done by a human authored script. Future work could see this
  be passed to the LLM as well. 
\item Uses Lua's {\tt loadstring()} to compile the code generated into a
  function {\tt run()}
\item Adds {\tt coalesce} as a key to the table {\tt lQDF}, the value of which
  is the function {\tt run} created above
  \ee

  \subsection{Calling an operator}
  \begin{verbatim}
local x = lQDF({ 1, 2, 3, }
local y = lQDF({ 4, 5, 6, })
local z = lQDF.coalesce(x, y)
  \end{verbatim}

This creates {\tt z}, an object of type {\tt lQDF}. 
Let's see take a peek under the covers. 
\be
\item Based on the types of the input arguments, {\tt x, y}, 
  determine whether a specialized function for the job exists. Let's assume it
  doesn't
\item Make substitutions in the \LaTeX\ file e.g., this is where we 
  would find that the {\tt qtype} of {\tt x,y} is {\tt FP64}, as 
  a result of which we would need to create a function whose name is 
  \verb+coalesece_FP64_FP64+
\item Assemble a larger \LaTeX\ file which contains 
  context which does not change across operators. 
  This is the information in Sections~\ref{data_format_defn} and 
  ~\ref{accessor_defn}.
\item Use {\tt pdflatex} to create a PDF with detailed instructions for the LLM.
  See Section~\ref{YYY} for an example.
\item POST the PDF to the LLM and get back 
  \be
\item the function declaration {\tt .h} file 
\item the function definition {\tt .c} file 
\item commands needed to create a {\tt .o} file and a {\tt .so} file. 
 Commands specify information such as
  \be
\item compiler to be used e..g, {\tt gcc} or {\tt ispc}
\item include paths
\item libraries that may need to be linked
\item compiler/linker flags 
  \ee
  \ee
\item Execute the commands above
\item \verb+ffi.cdef+ the {\tt .h} file and \verb+ffi.load+ the {\tt .so} file
  created above
\item The previous steps are skipped over if they had been executed before. 
The system remembers past work. Coming across a new symbol which necessitates
the above work can be thought of as a cache miss. 
This is an acceptable tradeoff {\bf if}
\be
\item these cache misses are infrequent 
\item the system is not sensitive to the latency they introduce,
  \ee
\item Provide C access to the Lua objects {\tt x, y, z}
\item Execute the C function using the pointers provided above as arguments
\item Return the Lua objects specified as the {\tt out\_args} in
  Section~\ref{XXX}.
  \ee

  \section{Conclusion}
  \subsection{Future Work}

\appendix
\section{Coalesce specification}
\label{coalesce_spec}
\verbatiminput{doc_coalesce.lua}

\input{doc_qdf_definition}
\bibliographystyle{alpha}
\bibliography{ref}

\end{document}

